% Chapter 4: Implementation
\chapter{Implementation}
\label{ch:implementation}

This chapter describes the development methodology and the key implementation
details of OrderVista, including algorithmic components, the AI integration
layer, and the offline synchronisation mechanism.

% ─────────────────────────────────────────────────────────────
\section{Development Methodology}
\label{sec:methodology}

An iterative, incremental approach was adopted over five phases spanning six
months (October 2025 to March 2026):

\begin{enumerate}
  \item \textbf{Foundation} (Oct) — project scaffolding, Supabase schema
        migration, authentication flow, basic CRUD screens.
  \item \textbf{Core Implementation} (Nov) — admin task management, CSV import
        with validation, driver task list with real-time subscriptions.
  \item \textbf{Feature Development} (Dec--Jan) — route generation with
        geocoding, proof-of-delivery workflow, AI summarisation and chat.
  \item \textbf{Integration} (Feb) — agentic AI actions, driver
        recommendation, fleet-wide optimisation, chat messaging.
  \item \textbf{Refinement} (Mar) — dynamic priority, mid-shift
        re-sequencing, performance metrics, security hardening, release build.
\end{enumerate}

Each phase concluded with a working APK build and a supervisor progress
review, ensuring that feedback could be incorporated before the next phase.
Version control followed a branch-per-feature model on GitHub with a single
integration branch (\texttt{feature/\allowbreak delivery-optimization-mvp}).

Figure~\ref{fig:mvvm} illustrates the layered MVVM architecture used
throughout the codebase.  Each layer communicates only with its immediate
neighbour, enforcing separation of concerns and testability.

\begin{figure}[ht]
  \centering
  \resizebox{0.82\textwidth}{!}{%
  \begin{tikzpicture}[
    box/.style={draw, rounded corners=4pt, minimum width=3cm,
                minimum height=0.9cm, font=\sffamily\small,
                align=center, thick},
    arr/.style={-{Stealth[length=3mm]}, thick, gray!70},
    node distance=1.4cm and 2.8cm
  ]
    % View layer
    \node[box, fill=teal!15]                     (v1) {Admin Portal\\Screens};
    \node[box, fill=teal!15, right=of v1]        (v2) {Driver App\\Screens};
    \node[box, fill=teal!15, right=of v2]        (v3) {Customer\\Tracking};

    % ViewModel layer
    \node[box, fill=orange!15, below=of v2]      (vm) {ViewModels\\{\footnotesize(TaskVM, RouteVM, ChatVM)}};

    % Repository layer
    \node[box, fill=blue!12, below=of vm]        (repo) {Repositories\\{\footnotesize(TaskRepo, RouteRepo)}};

    % Service layer
    \node[box, fill=violet!12, below left=1.4cm and 0.4cm of repo]  (svc1) {TaskService\\{\footnotesize(Supabase API)}};
    \node[box, fill=violet!12, below right=1.4cm and 0.4cm of repo] (svc2) {AIService\\{\footnotesize(DeepSeek R1)}};

    % Data layer
    \node[box, fill=green!12, below=1.4cm of svc1]  (drift) {Drift\\{\footnotesize(SQLite cache)}};
    \node[box, fill=green!12, below=1.4cm of svc2]  (supa)  {Supabase\\{\footnotesize(PostgreSQL)}};

    % Arrows
    \draw[arr] (v1.south) -- (vm);
    \draw[arr] (v2)       -- (vm);
    \draw[arr] (v3.south) -- (vm);
    \draw[arr] (vm)       -- (repo);
    \draw[arr] (repo)     -- (svc1);
    \draw[arr] (repo)     -- (svc2);
    \draw[arr] (svc1)     -- (drift);
    \draw[arr] (svc1)     -- (supa);
    \draw[arr] (svc2)     -- (supa);
  \end{tikzpicture}}%
  \caption{MVVM layered architecture: Views $\to$ ViewModels $\to$ Repositories $\to$ Services $\to$ Data.}
  \label{fig:mvvm}
\end{figure}

Figure~\ref{fig:class-diagram} expands the MVVM layers into a UML class
diagram showing the principal classes, their attributes, and
inter-class dependencies.

\begin{figure}[ht]
  \centering
  \resizebox{\textwidth}{!}{%
  \begin{tikzpicture}[
    class/.style={draw, rectangle split, rectangle split parts=3,
                  rectangle split part fill={#1!20, white, white},
                  rounded corners=2pt, font=\sffamily\scriptsize,
                  minimum width=3.8cm, thick, align=left},
    arr/.style={-{Stealth[length=2.5mm]}, thick, gray!70},
    comp/.style={-{Diamond[length=3mm, fill=black]}, thick, gray!70},
    uses/.style={-{Stealth[length=2.5mm]}, thick, dashed, gray!60},
    node distance=1.6cm and 1.4cm
  ]
    %% ── View layer ──────────────────────────────────
    \node[class=teal] (admin) {
      \textbf{AdminPortalScreen}
      \nodepart{second}
      \texttt{- \_tasks : List<DeliveryTask>}\\
      \texttt{- \_routes : List<Map>}\\
      \texttt{- \_selectedDriverId : String?}
      \nodepart{third}
      \texttt{+ \_runFleetOptimisation()}\\
      \texttt{+ \_buildDriverRecommendation()}\\
      \texttt{+ \_buildCustomerTracking()}\\
      \texttt{+ \_importCsv()}
    };

    \node[class=teal, right=2.2cm of admin] (tasklist) {
      \textbf{TaskListScreen}
      \nodepart{second}
      \texttt{- \_scannerKey : GlobalKey}
      \nodepart{third}
      \texttt{+ \_onBarcodeScanned()}\\
      \texttt{+ \_capturePoD()}\\
      \texttt{+ \_buildTaskCard()}
    };

    %% ── ViewModel layer ─────────────────────────────
    \node[class=orange, below=2cm of $(admin.south)!0.5!(tasklist.south)$] (taskvm) {
      \textbf{TaskViewModel}
      \nodepart{second}
      \texttt{- \_tasks : List<DeliveryTask>}\\
      \texttt{- \_isSyncing : bool}\\
      \texttt{- \_currentLocation : LatLng?}
      \nodepart{third}
      \texttt{+ loadTasks()}\\
      \texttt{+ computeEffectivePriority()}\\
      \texttt{+ resequenceFromLocation()}\\
      \texttt{+ syncPendingChanges()}
    };

    %% ── Repository layer ────────────────────────────
    \node[class=blue, below=2cm of taskvm] (taskrepo) {
      \textbf{TaskRepository}
      \nodepart{second}
      \texttt{- \_remote : TaskService}\\
      \texttt{- \_local : OrderVistaDatabase}
      \nodepart{third}
      \texttt{+ fetchTasks() : Future<List>}\\
      \texttt{+ upsertLocal()}\\
      \texttt{+ pushToRemote()}\\
      \texttt{+ enqueueSync()}
    };

    %% ── Service layer ───────────────────────────────
    \node[class=violet, below left=2cm and 0.8cm of taskrepo] (tasksvc) {
      \textbf{TaskService}
      \nodepart{second}
      \texttt{- \_supabase : SupabaseClient}
      \nodepart{third}
      \texttt{+ fetchAllTasks()}\\
      \texttt{+ updateTaskStatus()}\\
      \texttt{+ uploadProofPhoto()}
    };

    \node[class=violet, below right=2cm and 0.8cm of taskrepo] (aisvc) {
      \textbf{AiService}
      \nodepart{second}
      \texttt{- \_apiKey : String}\\
      \texttt{- \_model : String}
      \nodepart{third}
      \texttt{+ summariseTasks()}\\
      \texttt{+ chat()}\\
      \texttt{+ getAgenticSuggestion()}
    };

    \node[class=violet, right=1.8cm of tasklist] (mapbox) {
      \textbf{MapboxService}
      \nodepart{second}
      \texttt{- \_token : String}
      \nodepart{third}
      \texttt{+ geocodeAddress()}\\
      \texttt{+ nearestNeighbourSequence()}\\
      \texttt{+ getStaticMapUrl()}
    };

    %% ── Model layer ─────────────────────────────────
    \node[class=green, below=2cm of tasksvc, xshift=-0.5cm] (dtask) {
      \textbf{DeliveryTask}
      \nodepart{second}
      \texttt{+ id : String}\\
      \texttt{+ status : TaskStatus}\\
      \texttt{+ address : String}\\
      \texttt{+ latitude : double?}\\
      \texttt{+ longitude : double?}\\
      \texttt{+ scheduledStart : DateTime?}\\
      \texttt{+ priority : int}
      \nodepart{third}
      \texttt{+ toJson() : Map}\\
      \texttt{+ fromJson() : DeliveryTask}
    };

    \node[class=green, below=2cm of aisvc, xshift=0.5cm] (aiact) {
      \textbf{AiAction}
      \nodepart{second}
      \texttt{+ suggestion : String}\\
      \texttt{+ action : String}\\
      \texttt{+ taskId : String?}
      \nodepart{third}
      \texttt{+ fromJson() : AiAction}
    };

    %% ── Data layer ──────────────────────────────────
    \node[class=green, below=5.5cm of taskrepo] (driftdb) {
      \textbf{OrderVistaDatabase}
      \nodepart{second}
      \texttt{- \_executor : QueryExecutor}
      \nodepart{third}
      \texttt{+ watchTasks() : Stream}\\
      \texttt{+ insertTask()}\\
      \texttt{+ dequeueSyncItem()}
    };

    %% ── Inner class ─────────────────────────────────
    \node[class=green, below=1.4cm of mapbox] (geopoint) {
      \textbf{GeoPoint}
      \nodepart{second}
      \texttt{+ lat : double}\\
      \texttt{+ lon : double}
      \nodepart{third}
      \texttt{+ distanceTo() : double}
    };

    %% ── Relationships ───────────────────────────────
    \draw[arr] (admin)    -- (taskvm);
    \draw[arr] (tasklist) -- (taskvm);
    \draw[arr] (taskvm)   -- (taskrepo);
    \draw[arr] (taskrepo) -- (tasksvc);
    \draw[arr] (taskrepo) -- (driftdb);
    \draw[uses] (admin)   -- (aisvc);
    \draw[uses] (admin)   -- (mapbox);
    \draw[uses] (tasklist) -- (mapbox);
    \draw[arr] (tasksvc)  -- (dtask);
    \draw[arr] (aisvc)    -- (aiact);
    \draw[comp] (mapbox)  -- (geopoint);
    \draw[uses] (taskvm)  -- (aisvc);
  \end{tikzpicture}}%
  \caption{UML class diagram showing principal classes across MVVM layers, their attributes, operations, and dependencies.}
  \label{fig:class-diagram}
\end{figure}

% ─────────────────────────────────────────────────────────────
\section{Route Generation Pipeline}
\label{sec:route-impl}

Route generation is the most computationally significant feature.
Figure~\ref{fig:route-pipeline} shows the pipeline as an activity diagram, and
Algorithm~\ref{alg:nn} formalises the nearest-neighbour heuristic.

\begin{figure}[ht]
  \centering
  \resizebox{0.88\textwidth}{!}{%
  \begin{tikzpicture}[
    act/.style={draw, rounded corners=4pt, minimum width=3cm,
                minimum height=0.7cm, font=\sffamily\footnotesize,
                fill=blue!8, align=center},
    dec/.style={draw, diamond, aspect=2, inner sep=1pt,
                font=\sffamily\footnotesize, fill=orange!12},
    arr/.style={-{Stealth[length=2.5mm]}, thick},
    node distance=0.9cm
  ]
    \node[act]                     (start)  {Admin selects pending tasks};
    \node[act, below=of start]     (geo)    {Geocode each address\\(Mapbox API)};
    \node[dec, below=of geo]       (check)  {All geocoded?};
    \node[act, right=2cm of check] (skip)   {Flag failures for\\manual correction};
    \node[act, below=of check]     (nn)     {Run nearest-neighbour\\sequencer (Alg.~\ref{alg:nn})};
    \node[act, below=of nn]        (store)  {Persist route +\\route\_tasks};
    \node[act, below=of store]     (map)    {Generate static\\map preview};
    \node[act, below=of map]       (assign) {Admin confirms\\driver assignment};

    \draw[arr] (start)  -- (geo);
    \draw[arr] (geo)    -- (check);
    \draw[arr] (check)  -- node[left, font=\footnotesize]{Yes} (nn);
    \draw[arr] (check)  -- node[above, font=\footnotesize]{Partial} (skip);
    \draw[arr] (skip.south) |- (nn);
    \draw[arr] (nn)     -- (store);
    \draw[arr] (store)  -- (map);
    \draw[arr] (map)    -- (assign);
  \end{tikzpicture}}%
  \caption{Activity diagram for the route generation pipeline.}
  \label{fig:route-pipeline}
\end{figure}

\begin{algorithm}[ht]
  \caption{Greedy Nearest-Neighbour Route Sequencing}\label{alg:nn}
  \KwIn{Depot location $\mathbf{o}$; set of task locations $T = \{t_1, \dots, t_n\}$}
  \KwOut{Ordered sequence $S$}
  $S \leftarrow [\,]$\;
  $\mathit{current} \leftarrow \mathbf{o}$\;
  $R \leftarrow T$ \tcp*{remaining tasks}
  \While{$R \neq \emptyset$}{
    $t^* \leftarrow \arg\min_{t \in R}\; \textsc{Haversine}(\mathit{current},\; t)$\;
    Append $t^*$ to $S$\;
    $\mathit{current} \leftarrow t^*$\;
    $R \leftarrow R \setminus \{t^*\}$\;
  }
  \Return{$S$}\;
\end{algorithm}

The Haversine formula used in Algorithm~\ref{alg:nn} computes great-circle
distance:
\begin{equation}
  d = 2\,R_{\oplus}\;\arctan\!\Bigl(
    \sqrt{a},\; \sqrt{1-a}\Bigr),
  \quad
  a = \sin^{2}\!\frac{\Delta\phi}{2}
    + \cos\phi_1\,\cos\phi_2\;\sin^{2}\!\frac{\Delta\lambda}{2}
  \label{eq:haversine}
\end{equation}
where $R_{\oplus} = 6{,}371$\,km is the mean Earth radius, $\phi$ is
latitude, and $\lambda$ is longitude.  The implementation
(Listing~\ref{lst:haversine}) mirrors Equation~\ref{eq:haversine} directly.

\begin{lstlisting}[
  language=Java,
  caption={Haversine distance in Dart (\texttt{mapbox\_service.dart}).},
  label={lst:haversine},
  morekeywords={double, const, final, return}
]
double distanceTo(GeoPoint other) {
  const earthRadius = 6371.0; // km
  final dLat = _toRad(other.lat - lat);
  final dLng = _toRad(other.lng - lng);
  final a = sin(dLat/2) * sin(dLat/2) +
      cos(_toRad(lat)) * cos(_toRad(other.lat)) *
      sin(dLng/2) * sin(dLng/2);
  final c = 2 * atan2(sqrt(a), sqrt(1 - a));
  return earthRadius * c;
}
\end{lstlisting}

\paragraph{Complexity analysis.}
The greedy nearest-neighbour heuristic runs in $O(n^2)$ time, where $n$ is
the number of tasks.  While metaheuristics such as Ant Colony Optimisation
\autocite{dorigo1997ant} can yield shorter routes, the $O(n^2)$ cost is
acceptable for the typical task sets in last-mile delivery ($n < 50$) and,
critically, permits real-time mid-shift re-sequencing on a mobile device
without perceptible latency.

% ─────────────────────────────────────────────────────────────
\section{Dynamic Priority Recalculation}
\label{sec:priority-impl}

Task ordering is not solely distance-based.  A time-urgency bonus is added to
each task's base priority at sort time, ensuring that approaching time windows
are respected even when a closer but non-urgent task exists.

\begin{algorithm}[ht]
  \caption{Dynamic Priority with Time-Urgency Bonus}\label{alg:priority}
  \KwIn{Task $t$ with base priority $p$ and optional scheduled time $s$; current time $\mathit{now}$}
  \KwOut{Effective priority $p'$}
  \uIf{$t$ is delivered}{
    \Return{$p$}\;
  }
  \uIf{$s$ is null}{
    \Return{$p$}\;
  }
  $\Delta \leftarrow s - \mathit{now}$\;
  \uIf{$\Delta < 0$}{$\mathit{bonus} \leftarrow 50$ \tcp*{overdue}}
  \uElseIf{$\Delta \leq 30\,\text{min}$}{$\mathit{bonus} \leftarrow 30$ \tcp*{imminent}}
  \uElseIf{$\Delta \leq 60\,\text{min}$}{$\mathit{bonus} \leftarrow 20$}
  \uElseIf{$\Delta \leq 120\,\text{min}$}{$\mathit{bonus} \leftarrow 10$}
  \Else{$\mathit{bonus} \leftarrow 0$}
  \Return{$p + \mathit{bonus}$}\;
\end{algorithm}

The composite sort order used by the task list view model is:
(1)~status class (pending $>$ in-transit $>$ delivered),
(2)~effective priority (descending),
(3)~scheduled time (ascending),
(4)~stable tie-break by task ID.
This four-level comparator prevents task list flickering while still surfacing
urgent deliveries.

% ─────────────────────────────────────────────────────────────
\section{Driver Recommendation}
\label{sec:driver-rec}

When an admin generates a route, the system recommends the most suitable driver
using a composite score computed in a PostgreSQL function:

\begin{equation}
  \mathit{score} =
    \underbrace{w_a}_{\text{availability}}
    \times
    \underbrace{\frac{1}{1 + d_{\text{km}}}}_{\text{proximity}}
    \times
    \underbrace{\frac{1}{1 + n_{\text{tasks}}}}_{\text{workload}}
    \times
    \underbrace{\bigl(0.5 + 0.5 \cdot r_{\text{completion}}\bigr)}_{\text{performance}}
  \label{eq:rec-score}
\end{equation}

where $w_a \in \{0.5, 1.0\}$ is the availability weight, $d_{\text{km}}$ is
the PostGIS-computed distance from the driver's last known location to the task
cluster centroid, $n_{\text{tasks}}$ is the driver's current task count, and
$r_{\text{completion}}$ is the ratio of successfully delivered tasks to total
assigned tasks (historical).  The function is executed entirely within
PostgreSQL, leveraging PostGIS \texttt{ST\_Distance} for spatial computation
and avoiding a network round-trip for each candidate driver.

\paragraph{Alternatives considered.}
A client-side scoring approach was prototyped first but required fetching all
driver locations to the admin device.  Moving the computation server-side
reduced payload size and exploited PostGIS spatial indexes for sub-millisecond
distance queries.

% ─────────────────────────────────────────────────────────────
\section{Fleet-Wide Optimisation}
\label{sec:fleet-impl}

For dispatching multiple tasks across several drivers simultaneously, a greedy
nearest-driver assignment strategy is employed:

\begin{enumerate}
  \item Geocode all unassigned tasks.
  \item For each task, compute the Haversine distance to every available driver.
  \item Assign the task to the closest driver.
\end{enumerate}

This produces an $O(t \times d)$ assignment where $t$ is the number of tasks
and $d$ is the number of drivers.  While a Hungarian-algorithm assignment
would be globally optimal, the greedy approach was chosen because it runs in
real-time and the marginal optimality loss is negligible for small fleets
($d < 20$) typical of Hong Kong SMEs \autocite{toth2002vehicle}.

% ─────────────────────────────────────────────────────────────
\section{AI Integration}
\label{sec:ai-impl}

The AI subsystem comprises three modes of operation, each calling DeepSeek~R1
through the OpenRouter API \autocite{openrouter2026docs}.

\subsection{Task Summarisation}
\label{sec:ai-summary}

A compact system prompt instructs the model to summarise the driver's current
task list in under 160 tokens.  Temperature is set to $0.3$ to favour
deterministic, factual output.

\subsection{Free-Form Chat}
\label{sec:ai-chat}

A multi-turn chat interface injects the fleet context (task counts, route
status) as a system prompt, followed by the conversation history.  The model
responds with natural-language advice; tokens are capped at 200 to prevent
verbose answers.

\subsection{Agentic Suggestions}
\label{sec:ai-agentic}

The agentic mode extends the chat interface with a structured JSON action
schema.  The LLM is instructed to analyse the driver's situation and return a
JSON object containing a natural-language \texttt{suggestion} and an
executable \texttt{action} field from a predefined vocabulary:

\begin{center}
\small
\texttt{start\_next\_pending} \quad
\texttt{mark\_delivered} \quad
\texttt{go\_available} \quad
\texttt{go\_offline} \quad
\texttt{sync\_pending} \quad
\texttt{refresh\_tasks}
\end{center}

The client parses the JSON response and renders a one-tap confirmation chip.
Figure~\ref{fig:agentic-flow} shows the interaction sequence.

\begin{figure}[ht]
  \centering
  \resizebox{0.88\textwidth}{!}{%
  \begin{tikzpicture}[
    act/.style={draw, rounded corners=4pt, minimum width=2.8cm,
                minimum height=0.65cm, font=\sffamily\footnotesize,
                fill=blue!8, align=center},
    dec/.style={draw, diamond, aspect=2.2, inner sep=1pt,
                font=\sffamily\footnotesize, fill=orange!12, align=center},
    arr/.style={-{Stealth[length=2.5mm]}, thick},
    node distance=0.85cm
  ]
    \node[act]                     (ctx)   {Build context\\{(}tasks, status, time{)}};
    \node[act, below=of ctx]       (call)  {Call DeepSeek R1\\(JSON action schema)};
    \node[dec, below=of call]      (parse) {Valid JSON?};
    \node[act, below=of parse]     (chip)  {Render suggestion +\\action chip};
    \node[act, right=2.2cm of parse] (fallback) {Show plain-text\\advisory};
    \node[dec, below=of chip]      (tap)   {User confirms?};
    \node[act, below=of tap]       (exec)  {Execute action\\(local state change)};
    \node[act, right=2.2cm of tap] (ignore) {No-op};

    \draw[arr] (ctx)      -- (call);
    \draw[arr] (call)     -- (parse);
    \draw[arr] (parse)    -- node[left, font=\footnotesize]{Yes} (chip);
    \draw[arr] (parse)    -- node[above, font=\footnotesize]{No}  (fallback);
    \draw[arr] (chip)     -- (tap);
    \draw[arr] (tap)      -- node[left, font=\footnotesize]{Yes} (exec);
    \draw[arr] (tap)      -- node[above, font=\footnotesize]{No}  (ignore);
  \end{tikzpicture}}%
  \caption{Agentic AI suggestion flow with user confirmation gate.}
  \label{fig:agentic-flow}
\end{figure}

This human-in-the-loop design prevents the LLM from executing unintended
actions—a critical safeguard given that LLM outputs can hallucinate
\autocite{ji2024hallucination}.  Temperature is lowered to $0.2$ for the
agentic mode to maximise output predictability.

% ─────────────────────────────────────────────────────────────
\section{Proof-of-Delivery Workflow}
\label{sec:pod-impl}

Before a task can be marked as \emph{completed}, the driver must capture a
photo using the device camera.  The image is written to Supabase Storage with a
path derived from the task ID, and the public URL is written back to the
\texttt{tasks.proof\_photo\_url} column.  The admin portal displays the photo
in a review panel alongside the delivery address and completion timestamp.

An optional confirmation code (e.g.\ a PIN provided by the recipient) is
stored in \texttt{tasks.proof\_code}.  Both fields together form the proof
of delivery (PoD) bundle, which is auditable per task.

% ─────────────────────────────────────────────────────────────
\section{Offline Synchronisation}
\label{sec:sync-impl}

The offline sync mechanism (Figure~\ref{fig:sync}) ensures data durability in
areas with unreliable connectivity.

\begin{figure}[ht]
  \centering
  \resizebox{0.92\textwidth}{!}{%
  \begin{tikzpicture}[
    act/.style={draw, rounded corners=4pt, minimum width=3cm,
                minimum height=0.65cm, font=\sffamily\footnotesize,
                fill=blue!8, align=center},
    store/.style={draw, cylinder, shape border rotate=90,
                  minimum width=1.6cm, minimum height=0.8cm,
                  font=\sffamily\footnotesize, fill=green!10, align=center},
    arr/.style={-{Stealth[length=2.5mm]}, thick},
    darr/.style={-{Stealth[length=2.5mm]}, thick, dashed, red!60},
    node distance=1cm and 2cm
  ]
    \node[act]                         (write)  {Driver action\\(status update, PoD)};
    \node[store, right=of write]       (drift)  {Drift\\{(}SQLite{)}};
    \node[act, right=of drift]         (queue)  {Sync queue\\{(}exponential backoff{)}};
    \node[store, right=of queue]       (supa)   {Supabase\\{(}PostgreSQL{)}};

    \draw[arr]  (write) -- (drift);
    \draw[arr]  (drift) -- (queue);
    \draw[arr]  (queue) -- node[above, font=\footnotesize]{Online} (supa);
    \draw[darr] (queue) to[out=-30,in=-150] node[below, font=\footnotesize]{Retry [0,\,10,\,30,\,60,\,120]\,s} (queue);

    % Realtime path back
    \draw[arr, blue!60] (supa) to[out=150,in=30]
      node[above, font=\footnotesize, blue!60]{Realtime subscription} (drift);
  \end{tikzpicture}}%
  \caption{Offline synchronisation with exponential-backoff retry and Realtime reconciliation.}
  \label{fig:sync}
\end{figure}

\begin{enumerate}
  \item Every write first persists to the local Drift database (immediate UI
        feedback).
  \item The mutation is enqueued into \texttt{sync\_queue\_items} with
        \texttt{status = 'pending'}.
  \item A background isolate dequeues items and attempts upload.  On failure,
        \texttt{attempts} is incremented and \texttt{last\_attempt\_at} is
        updated; the next retry is delayed according to a progressive
        schedule $[0, 10, 30, 60, 120]$~seconds.  After five failed attempts
        the item is marked as permanently failed.
  \item On reconnection, Supabase Realtime pushes server-side changes;
        conflicts are resolved by comparing \texttt{updated\_at} timestamps
        (last-writer-wins).
\end{enumerate}

% ─────────────────────────────────────────────────────────────
\section{Admin Portal}
\label{sec:admin-impl}

The admin portal is implemented as a role-gated view within the same Flutter
application.  Key screens include:

\begin{itemize}
  \item \textbf{Task Management} — CRUD operations with inline validation;
        bulk CSV import parses headers (\texttt{address}, \texttt{notes},
        \texttt{priority}, \texttt{scheduled\_start}, \texttt{scheduled\_end})
        and rejects malformed rows with error feedback.
  \item \textbf{Route Dashboard} — displays generated routes on a Mapbox
        static-map tile with sequenced stop markers; shows total distance and
        recommended driver.
  \item \textbf{Fleet Monitor} — real-time driver locations plotted on a map
        via Supabase Realtime subscriptions to the \texttt{company\_members}
        table (which stores \texttt{last\_known\_location}).
  \item \textbf{PoD Review} — photo and signature viewer per task, with
        accept/reject controls.
  \item \textbf{Chat} — admin-to-driver messaging with broadcast support.
\end{itemize}

% ─────────────────────────────────────────────────────────────
\section{Mid-Shift Re-Sequencing}
\label{sec:reseq-impl}

When conditions change after a route has been dispatched (e.g.\ a task is
cancelled or traffic reroutes the driver), the system can re-sequence the
remaining tasks from the driver's \emph{current GPS position} rather than the
original depot.  The existing nearest-neighbour sequencer
(Algorithm~\ref{alg:nn}) is re-invoked with the driver's live coordinates as
the origin, and only undelivered tasks are included.  The result is stored as
a temporary client-side ordering override that expires after one hour,
preventing stale overrides from persisting across shifts.

% ─────────────────────────────────────────────────────────────
\section{Barcode and QR-Code Scanner}
\label{sec:barcode-impl}

Drivers can scan a barcode or QR code affixed to a parcel to instantly locate
the corresponding task in their list.  The feature is implemented as a
\texttt{BarcodeScannerSheet}—a reusable bottom-sheet widget backed by the
\texttt{mobile\_scanner} package, which accesses the device camera and
performs real-time barcode detection.

The interaction proceeds as follows:

\begin{enumerate}
  \item The driver taps the scan icon on the task list toolbar.
  \item A bottom sheet opens containing a live camera preview
        (\texttt{MobileScanner} widget).
  \item On detection, a debounce flag (\texttt{\_hasResult}) prevents
        duplicate scans; the first valid code is trimmed and returned via
        \texttt{Navigator.pop(value)}.
  \item The calling screen receives the scanned string and filters the task
        list to the matching delivery.
\end{enumerate}

This eliminates the need for manual task lookup when parcels carry
machine-readable labels—a common workflow in Hong Kong logistics depots
where hundreds of parcels are handled per shift.

% ─────────────────────────────────────────────────────────────
\section{Customer Tracking Lookup}
\label{sec:tracking-impl}

The admin portal includes a customer-facing tracking lookup that allows
administrators—or, in a future extension, end-customers directly—to query
delivery status using a human-readable tracking code (format
\texttt{OV-ABC1-XYZ2}).

\subsection{Implementation}

The lookup is implemented via a Supabase RPC call to the
\texttt{lookup\_delivery\_by\_tracking} PostgreSQL function, which accepts a
single parameter \texttt{p\_tracking\_code} and returns a JSON object
containing:

\begin{itemize}
  \item \texttt{found} — boolean indicating whether a matching delivery exists,
  \item \texttt{status} — current task state (pending / in-transit / delivered),
  \item \texttt{address} — delivery address,
  \item \texttt{driver\_assigned} — name of the assigned driver,
  \item \texttt{proof\_available} — boolean indicating whether PoD has been captured,
  \item \texttt{last\_update} — ISO-8601 timestamp of the most recent status change.
\end{itemize}

The UI renders the result in a colour-coded card: green for delivered, blue for
in-transit, and orange for pending.  Error states (no match, network failure)
are presented as inline messages, keeping the admin within the same screen
context.
